\documentclass[11pt,a4paper]{article}
\usepackage[margin=1in]{geometry}
\usepackage{hyperref}
\usepackage{enumitem}
\usepackage{graphicx}
\usepackage{xcolor}

% -----------------------------------------------------
% bvraghav-tiet-ucs503p-article.sty
% -----------------------------------------------------

\makeatletter \newcommand{\BvrSetLabInstructor}[1]{%
  \gdef\Bvr@LabInstructor{#1}}
\newcommand{\Bvr@LabInstructor}{Dr. Raghav %
  B. Venkataramaiyer}
\newcommand{\BvrLabInstructorValue}{\Bvr@LabInstructor}
\makeatother

\usepackage{titling}
\pretitle{\begin{center}\bfseries\fontsize{18bp}{18bp}\selectfont}
\makeatletter
\newcommand{\BvrSubtitle}[1]{%
  \posttitle{%
    \par\end{center}
    \begin{center}\large#1\end{center}
    \vskip 0.5em}%
}
\makeatother

\newcommand{\BvrHint}[1]{%
  {\normalfont\normalsize\itshape\hfill #1}}

% -----------------------------------------------------

\title{SportSathi: An AI-Based Fitness Screening and Sports Talent Guidance Platform for Schools}

\BvrSetLabInstructor{Dr. Raghav B. Venkataramaiyer}

\BvrSubtitle{%
  Project Proposal for UCS503P  \\
  Submitted to: \BvrLabInstructorValue \\ \bfseries
  Thapar Institute of Engineering and Technology%
}

\author{
  Kiratpreet Kaur, CSED%
  \thanks{Roll No: \texttt{1024030402}, Batch: \texttt{3C31}, %
    Email: \texttt{kkaur1\_be24\@thapar.edu}}
  \and
  Prisha Goyal, CSED%
  \thanks{Roll No: \texttt{1024030404}, Batch: \texttt{3C31}, %
    Email: \texttt{pgoyal1\_be24\@thapar.edu}}
}

\date{\today}

\begin{document}
\maketitle

\section{Higher-order goal%
  \BvrHint{\ldots secondary framing}}
This project aims to help schools move away from
ad-hoc, paper-based fitness records and toward
evidence-informed sports guidance, so that students with
genuine athletic potential are less likely to be
overlooked and students who are interested in sport
receive individualised, data-backed direction rather than
generic advice. The underlying goal of producing more
capable, confident student-athletes is broader than any
single school or sport; the immediate engineering
objective is narrower: build a measurable, deployable
fitness-screening and sport-suitability guidance platform
that keeps a coach in control of every recommendation.

\section{Time-to-value %
\BvrHint{\ldots primary heuristic}}
\begin{itemize}[leftmargin=0.7cm]
    \item Meaningful increments will be delivered quickly by prototyping the core assessment-to-prediction workflow first -- a coach records a student's general fitness scores and receives a machine-learning-generated fitness-level prediction -- before layering on sport-specific tracking and progression forecasting.
    \item Fast feedback is prioritised over perfection: the classification model will be validated in isolation on an existing student fitness dataset before integration, so that data-quality problems surface within days rather than at the end of the semester.
    \item Success in the first iteration is defined by what can be shown and measured: a coach can record an assessment and receive an explainable, confidence-scored fitness prediction, gated behind mandatory coach approval before the student or parent ever sees it. Sport-specific recommendation logic and long-term progression tracking are planned for later iterations.
\end{itemize}

\section{Problem Statement}
Schools conduct sports activities regularly, but student
fitness and performance data is rarely recorded in a
structured, comparable way. Students with genuine athletic
potential can go unnoticed simply because no one is
systematically comparing fitness data across the school.
Students who are interested in sport rarely receive
individualised guidance on which sport suits their
physical profile or what to improve, and without repeated,
comparable assessments, neither the student nor the coach
can reliably tell whether training is working. Coaches
lack a structured way to compare students across age
groups, fitness levels, or improvement needs.

This matters now because a standard, low-cost fitness-test
battery (in the spirit of the Khelo India protocol)
already gives schools a consistent way to measure
students; the missing piece is a system that turns
repeated test results into structured, coach-reviewed
guidance rather than a one-time paper record.

\section{Proposed Solution %
\BvrHint{\ldots Software Proposition}}
\subsection{Overview}
SportSathi will be a web-based fitness-screening and
sports-guidance platform built around the following
components:
\begin{itemize}[leftmargin=0.7cm]
    \item \textbf{Fitness assessment module:} records a coach-entered general fitness assessment for a student, covering strength, endurance, flexibility, speed-agility, BMI, attendance, and motivation.
    \item \textbf{ML-assisted fitness classification:} predicts a student's fitness category and a confidence score from the recorded assessment, using an explainable classification model.
    \item \textbf{Sport suitability recommendation:} recommends 2--3 suitable sports for the student's fitness profile, with a stated, plain-language reason.
    \item \textbf{Coach approval workflow:} holds every prediction back from students and parents until a coach explicitly approves or overrides it.
    \item \textbf{Student and parent guidance views:} display approved results, together with historical assessments and progress charts.
    \item \textbf{PDF report generation:} produces a downloadable summary report per student.
\end{itemize}

\subsection{Core Workflow}
A coach will record a student's general fitness assessment
through the web frontend. The backend validates and stores
the assessment, then invokes the machine learning service,
which returns a predicted fitness level, a confidence
score, and 2--3 recommended sports with stated reasons. The
coach reviews the prediction on the coach dashboard and
marks it approved or overridden. Only approved results
become visible in the student and parent views, alongside
historical assessments and a progress comparison against
the student's own prior results.

\subsection{Deliberate Scope Narrowing: Two-Tier Design}
Collecting every conceivable metric for every possible
sport, for every student, from day one would make the
system too complex for a coach to realistically operate.
The data model is therefore split into two tiers:
\textbf{Tier 1 -- General Athleticism} is a single universal
fitness battery applied to every student and used to
recommend suitable sports; this is the complete scope of
the current semester. \textbf{Tier 2 -- Sport-Specific
KPIs} is a small, five-to-seven-indicator set applied only
to students funnelled into a specific sport squad, together
with a progression model comparing them to age-group
benchmarks; this tier is explicitly treated as future
scope.

\subsection{Operational Constraints for an Educational, School-Access Setting}
Keep the assessment-entry and dashboard interfaces usable
on low-to-mid range devices via responsive web design.
Avoid dependence on advanced or research-grade modelling;
favour explainable, inspectable models (Decision Tree,
Random Forest) over black-box architectures, since a coach
must be able to trust and explain a recommendation to a
parent. Design for deployability using common web hosting
and a managed relational database.

\section{Solution Approach %
\BvrHint{\ldots Engineering focus}}
This is an engineering course project, so the focus is on
correct data modelling, honest model evaluation, and
workflow automation rather than novel machine learning
research.

\subsection{Fitness Classification %
\BvrHint{\ldots non-research}}
Fitness-level prediction will use standard, explainable
models rather than research-grade architectures: a
Decision Tree and a Random Forest classifier, evaluated
against each other, with class-imbalance handling (class
weighting and SMOTE oversampling) investigated given the
expected rarity of a top-tier fitness category in real
school data. Before any accuracy figure is trusted,
features that could leak the target label will be
identified and excluded from the training input. The model
will be surfaced to the user only as decision support: a
coach must approve or override every prediction before it
is acted on.

\subsection{Sport Recommendation %
\BvrHint{\ldots non-research}}
Sport suitability will be a rule-based mapping from a
student's fitness profile to a small set of candidate
sports, rather than a learned recommendation model, keeping
the logic auditable and explainable to a coach or parent.

\subsection{Web Architecture %
\BvrHint{\ldots web-based solution}}
A typical three-tier web architecture will be used:
\begin{itemize}[leftmargin=0.7cm]
    \item \textbf{Frontend:} React/Vite interface for Student, Coach, Admin, and Parent roles.
    \item \textbf{Backend API:} Node.js/Express, handling authentication, role checks, validation, and orchestration of the ML service.
    \item \textbf{Machine learning service:} a Python/scikit-learn service invoked per assessment.
    \item \textbf{Database:} PostgreSQL, storing student profiles, assessments, coaches, recommendations, and administrative accounts in a normalised relational schema.
\end{itemize}

\subsection{CI/CD and Fast Delivery}
Automatically build and run backend unit tests and linting
on every change; produce a repeatable deployment artefact
for a staging environment; and enable safe, frequent
releases of incremental features, with the Tier 1
assessment-to-prediction workflow first, and sport
recommendation and reporting features after.

\section{Evaluation Criteria %
\BvrHint{\ldots measurable and attributable}}
Success will be defined using measurable metrics tied to
the core workflow.

\subsection{Primary Evaluation Metric}
\textbf{Classification quality on a held-out test set:}
precision, recall, and macro-averaged F1-score per fitness
category, with particular attention to per-class
performance rather than overall accuracy alone.

\textbf{Target:} an explainable classifier (Decision Tree
or Random Forest) that achieves at least 80\% overall
accuracy on a held-out split of the assessment dataset,
evaluated after excluding any leaky/derived input features.

\textbf{Attribution:} measured directly via
\texttt{scikit-learn}'s classification report and confusion
matrix on the trained pipeline.

\subsection{Secondary Evaluation Metrics}
\begin{itemize}[leftmargin=0.7cm]
    \item \textbf{Workflow integrity:} a prediction must be verified inaccessible to the student/parent view until a coach-approval action is recorded.
    \item \textbf{Role isolation:} each of the four roles (Student, Coach, Admin, Parent) must access only its intended views and data.
    \item \textbf{Recommendation relevance:} percentage of sport recommendations a reviewing coach agrees are plausible given the student's fitness profile.
    \item \textbf{Coach usability:} time taken by a coach to complete a full assessment entry, kept low enough to be practical for an entire class or cohort.
\end{itemize}

\subsection{Pilot Validation Plan}
Seed the system with a synthetic or openly available
student fitness dataset covering multiple age groups and
fitness dimensions. Recruit a small number of pilot
reviewers to play the Coach role, entering sample
assessments and reviewing/approving generated
recommendations during a short trial window. Compare the
model's fitness-level predictions against the dataset's
known labels, and gather simple clarity feedback on the
coach dashboard from pilot reviewers.

\section{Scalability %
\BvrHint{\ldots with foresight}}
\textbf{Scalable in theory:} the schema will be normalised
around one core fact -- a student is assessed repeatedly
over time -- so the natural unit of storage is one row per
assessment, not one row per student, supporting growth in
students, assessments, and concurrent usage without a
redesign. The schema will be deliberately additive, so a
future Tier 2 sport-specific metrics table can be
introduced without modifying the Tier 1 tables.

\textbf{Deployable in practice:} the frontend, backend, and
machine learning service will be decoupled and
independently deployable, and the machine learning service
will be invoked as a stateless per-request call, avoiding a
persistent, hard-to-scale process.

\section{Technology Stack}
A mature set of tools and services supports this project:
\begin{itemize}[leftmargin=0.7cm]
    \item \textbf{Frontend framework:} React / Vite.
    \item \textbf{Web framework for REST APIs:} Node.js / Express.
    \item \textbf{Relational database:} PostgreSQL.
    \item \textbf{Machine learning:} Python, scikit-learn, imbalanced-learn.
    \item \textbf{Authentication approach:} JWT, token-based.
    \item \textbf{Data source:} a public/open student physical-education performance dataset, supplemented with synthetic records generated to Khelo India-style benchmark ranges.
    \item \textbf{Test framework:} Pytest / Jest for backend and model-evaluation testing.
    \item \textbf{CI runner and staging target:} GitHub Actions with a free-tier host.
\end{itemize}
These mature components let the team focus on data
modelling, workflow correctness, and honest model
evaluation rather than inventing new infrastructure.

\section{Project Scope and Deliverables %
\BvrHint{\ldots incremental delivery}}
\subsection{Initial Deliverable (First Iteration)}
\begin{itemize}[leftmargin=0.7cm]
    \item Authentication and role-based access (Student / Coach / Admin / Parent).
    \item Tier 1 fitness assessment entry with automatic BMI and improvement-percentage calculation.
    \item Trained and evaluated fitness-level classifier (leakage-checked, class-imbalance-aware).
    \item Coach approval/override workflow, enforced before student or parent visibility.
    \item CI: build, backend unit tests, and linting.
    \item CD to a staging environment with a single pipeline job.
\end{itemize}

\subsection{Subsequent Deliverables}
\begin{itemize}[leftmargin=0.7cm]
    \item Sport-suitability recommendation logic (fitness-profile-to-sport rule engine).
    \item Student and parent dashboards with progress charts and PDF report generation.
    \item Admin dashboard with school-wide distribution analytics.
    \item Tier 2 sport-specific KPI tracking for squad-selected students, added one sport at a time.
    \item A progression/forecasting model comparing a student's improvement trajectory against age-group benchmarks.
\end{itemize}
Full production-grade privacy/consent workflows,
multi-school deployment, and mobile applications are
acknowledged as part of the broader problem but are treated
as future scope beyond this course project.

\section{Risks and Mitigations}
\begin{itemize}[leftmargin=0.7cm]
    \item \textbf{Data leakage / inflated accuracy:} validate that no derived/summary column is included as a model input, and report per-class metrics rather than overall accuracy alone.
    \item \textbf{Class imbalance:} a rare fitness category may be too small to predict reliably; mitigate with class weighting and SMOTE, and document honestly as a limitation.
    \item \textbf{Coach adoption:} keep assessment entry to a single form with automatic BMI calculation, and never surface an unreviewed prediction to a student or parent.
    \item \textbf{Data privacy / consent:} use only synthetic or openly available fitness data during development; no real student-identifying data will be collected.
    \item \textbf{Scope creep from sport-specific tracking:} mitigated directly by the two-tier design, which confines the current semester to a single, universal fitness battery.
    \item \textbf{Operational issues in the pilot:} exercise the staging environment early and keep CI/CD producing repeatable, one-command deployments.
\end{itemize}

\section{Summary}
This project proposes SportSathi, a deployable,
coach-supervised web platform that will turn a school's
fitness assessment data into structured, explainable
sport-suitability guidance, without ever letting a machine
learning prediction reach a student or parent unreviewed.
It meets the course criteria by emphasising time-to-value
through incremental delivery, planning CI/CD to support
frequent and safe releases, and using measurable evaluation
criteria, especially per-class classification quality and
workflow integrity. The project stays focused on
engineering correctness, honest evaluation, and a
deliberately scoped two-tier data strategy, rather than on
research-driven modelling novelty.

\end{document}
